\chapter{พลาสมาเกษตรอัจฉริยะ ปัญญาประดิษฐ์ และแนวโน้มในอนาคต}
\label{chap:ch11}

% =========================================================================
% แผนบริหารการสอนประจำบทที่ 11 (Modular Lesson Plan Boxes - 2-Page Standard)
% =========================================================================

\vspace*{-6pt}
\section*{แผนบริหารการสอนประจำบทที่ 11}
\addcontentsline{toc}{section}{แผนบริหารการสอนประจำบทที่ 11}

% --- Page 1 of Lesson Plan ---
\begin{planbox}{หัวข้อเนื้อหาประจำบท}
\begin{enumerate}
    \item สถาปัตยกรรมระบบพลาสมากายภาพไซเบอร์สำหรับการเกษตรแม่นยำ
    \item เครือข่ายโทรมาตร IoT และการตรวจวัดทางเคมีกายภาพแบบพหุมิติ
    \item การควบคุมแบบวงรอบปิดด้วยการเรียนรู้แบบเสริมกำลัง (Reinforcement Learning)
    \item คอมพิวเตอร์วิทัศน์บนขอบสำหรับการประเมินฟีโนไทป์ของเมล็ดพืชและการตรวจจับโรคหลังการเก็บเกี่ยว
    \item แพลตฟอร์มพลาสมาพลังงานแสงอาทิตย์เพื่อการเกษตรคาร์บอนสุทธิเป็นศูนย์
    \item แผนที่นำทางเทคโนโลยีและการขยายผลสู่อุตสาหกรรมในอนาคต (พ.ศ. 2569--2578)
\end{enumerate}
\end{planbox}

\begin{planbox}{วัตถุประสงค์เชิงพฤติกรรม}
\begin{enumerate}
    \item อธิบายสถาปัตยกรรมของระบบพลาสมากายภาพไซเบอร์ (Cyber-Physical Plasma Systems) ได้อย่างเป็นระบบ
    \item ออกแบบผังการเชื่อมต่อเซนเซอร์โทรมาตร IoT ในการตรวจวัดพารามิเตอร์น้ำ PAW และการคายประจุได้
    \item ประยุกต์ใช้อัลกอริทึมการเรียนรู้แบบเสริมกำลัง (PPO/DQN) ในการควบคุมกำลังไฟฟ้าและอัตราส่วน RONS ได้
    \item บูรณาการโมเดลคอมพิวเตอร์วิทัศน์บนอุปกรณ์ขอบ (Edge AI) ในการคัดแยกเมล็ดพืชและตรวจจับเชื้อราได้
    \item วิเคราะห์ความคุ้มค่าเชิงพลังงานและผลกระทบต่อสิ่งแวดล้อมตามแนวทางเศรษฐกิจหมุนเวียนคาร์บอนต่ำได้
\end{enumerate}
\end{planbox}

\newpage

% --- Page 2 of Lesson Plan ---
\begin{planbox}{วิธีสอนและกิจกรรมการเรียนการสอน}
\begin{enumerate}
    \item บรรยายเชิงทัศนะควบคู่กับการนำเสนอกรณีศึกษาฟาร์มอัจฉริยะระบบพลาสมาอัตโนมัติในระดับนานาชาติ
    \item อภิปรายเชิงปฏิบัติการเกี่ยวกับการออกแบบฟังก์ชันการให้รางวัล (Reward Function) สำหรับเอเจนท์ควบคุมพลาสมา
    \item สาธิตการเชื่อมต่อไมโครคอนโทรลเลอร์ ESP32 และเซนเซอร์ pH/ORP ผ่านโพรโทคอล MQTT สู่แดชบอร์ดคลาวด์
    \item วิเคราะห์ภาพถ่ายการงอกของเมล็ดพืชด้วยโมเดลตรวจจับวัตถุ YOLO บนอุปกรณ์ประมวลผลขอบ
    \item วางแผนจัดทำข้อเสนอโครงการนวัตกรรมพลาสมาเพื่อสิ่งแวดล้อมที่สอดคล้องกับเป้าหมายการพัฒนาที่ยั่งยืน (SDGs)
\end{enumerate}
\end{planbox}

\begin{planbox}{สื่อการเรียนการสอน}
\begin{enumerate}
    \item ชุดทดสอบระบบพลาสมาอัตโนมัติควบคุมด้วยบอร์ดสมองกลฝังตัวและระบบขับเคลื่อนสเต็ปเปอร์มอเตอร์
    \item แพลตฟอร์มจำลองสภาพแวดล้อม Gym/Gymnasium สำหรับการฝึกสอน Reinforcement Learning
    \item ชุดข้อมูลภาพถ่ายความละเอียดสูงของเมล็ดข้าว ข้าวโพด และผลผลิตหลังการเก็บเกี่ยวที่ผ่านการกระตุ้นด้วยพลาสมา
    \item รายงานวิจัยและฐานข้อมูลสิทธิบัตรเทคโนโลยีพลาสมาการเกษตรสากล (ปี ค.ศ. 2024--2026)
\end{enumerate}
\end{planbox}

\begin{planbox}{การวัดและประเมินผล}
\begin{enumerate}
    \item ประเมินจากการนำเสนอแนวคิดการออกแบบระบบพลาสมากายภาพไซเบอร์ ร้อยละ 20
    \item ประเมินความถูกต้องของแบบจำลอง Reinforcement Learning ในการควบคุมความเข้มข้นสารละลาย ร้อยละ 20
    \item ประเมินผลงานการพัฒนาโมเดล Edge AI ตรวจจับคุณภาพเมล็ดพันธุ์ ร้อยละ 10
    \item ประเมินจากการสอบข้อเขียนประจำภาคการศึกษา ร้อยละ 50
\end{enumerate}
\end{planbox}

\clearpage

% =========================================================================
% เนื้อหาบทเรียน (เริ่มต้นหน้า 3 อย่างเป็นระเบียบ)
% =========================================================================

\section{สถาปัตยกรรมระบบพลาสมากายภาพไซเบอร์}

ในยุคเกษตรกรรมอัจฉริยะ (Smart Agriculture 5.0) การนำเทคโนโลยีพลาสมามาประยุกต์ใช้ไม่ได้จำกัดอยู่เพียงการเปิดเครื่องทำงานแบบตั้งค่าคงที่ (Static open-loop operation) แต่ได้ก้าวข้ามไปสู่ \textbf{ระบบพลาสมากายภาพไซเบอร์ (Cyber-Physical Plasma Systems หรือ CPPS)} ซึ่งผสานรวมระหว่างฟิสิกส์การคายประจุจริงในแปลงเพาะปลูก เซนเซอร์ตรวจวัดทางชีวกายภาพ เครือข่ายโทรมาตร IoT และอัลกอริทึมปัญญาประดิษฐ์บนคลาวด์และอุปกรณ์ขอบ

เป้าหมายสำคัญของ CPPS คือการบรรลุภาวะ \textbf{การบำบัดพลาสมาเฉพาะจุดและแม่นยำ (Precision Plasma Treatment)} โดยระบบจะปรับเปลี่ยนสภาวะการคายประจุ (แรงดันไฟฟ้า, ความถี่, ดิวตีไซเคิล, อัตราการไหลของก๊าซ) ให้สอดคล้องกับความต้องการของพืชและระดับการเข้าทำลายของเชื้อก่อโรคในแต่ละวินาที เพื่อให้ได้ประสิทธิภาพสูงสุดโดยไม่สูญเสียพลังงานและไม่ทำให้พืชเกิดความเครียดจากภาวะออกซิเดชันเกินขนาด (Oxidative stress)

\begin{figure}[H]
\centering
\begin{tikzpicture}[scale=0.78, every node/.style={transform shape}]
    % Physical Layer
    \draw[thick, fill=blue!10, rounded corners=6pt] (-0.5, 0) rectangle (10.5, 2.2);
    \node[anchor=west] at (0, 1.8) {\textbf{\color{blue!80!black}1. ชั้นกายภาพ (Physical Plasma Layer)}};
    \node[draw, fill=white, rounded corners=4pt] at (2.2, 0.9) {\small รีแอคเตอร์ DBD/APPJ};
    \node[draw, fill=white, rounded corners=4pt] at (5.5, 0.9) {\small แปลงปลูก / เมล็ด / น้ำ PAW};
    \node[draw, fill=white, rounded corners=4pt] at (8.8, 0.9) {\small หัววัด pH/ORP/OES};
    \draw[thick, ->] (3.6, 0.9) -- (4.2, 0.9);
    \draw[thick, ->] (6.8, 0.9) -- (7.5, 0.9);

    % Edge AI Layer
    \draw[thick, fill=teal!10, rounded corners=6pt] (-0.5, 3.0) rectangle (10.5, 5.2);
    \node[anchor=west] at (0, 4.8) {\textbf{\color{teal!80!black}2. ชั้นประมวลผลขอบ (Edge AI \& Telemetry Layer)}};
    \node[draw, fill=white, rounded corners=4pt] at (2.2, 3.9) {\small ขับเคลื่อนพัลส์นาโนวินาที};
    \node[draw, fill=white, rounded corners=4pt] at (5.5, 3.9) {\small ไมโครคอนโทรลเลอร์ ESP32};
    \node[draw, fill=white, rounded corners=4pt] at (8.8, 3.9) {\small กล้องสมาร์ทวิชัน Edge AI};
    \draw[thick, <-] (3.6, 3.9) -- (4.2, 3.9);
    \draw[thick, <-] (6.8, 3.9) -- (7.5, 3.9);

    % Cyber / Cloud Layer
    \draw[thick, fill=purple!10, rounded corners=6pt] (-0.5, 6.0) rectangle (10.5, 8.2);
    \node[anchor=west] at (0, 7.8) {\textbf{\color{purple!80!black}3. ชั้นไซเบอร์และสมองกลอัจฉริยะ (Cyber-Brain Layer)}};
    \node[draw, fill=white, rounded corners=4pt] at (3.0, 6.9) {\small โมเดล PINN \& RL Agent};
    \node[draw, fill=white, rounded corners=4pt] at (7.5, 6.9) {\small คลาวด์แดชบอร์ด \& ดิจิทัลทวิน};
    \draw[thick, <->] (4.8, 6.9) -- (5.7, 6.9);

    % Inter-layer arrows
    \draw[very thick, ->, blue!70!black] (8.8, 2.2) -- (8.8, 3.0) node[midway, right] {\scriptsize ข้อมูลเซนเซอร์ดิบ};
    \draw[very thick, ->, teal!70!black] (8.8, 5.2) -- (8.8, 6.0) node[midway, right] {\scriptsize ส่งผ่าน MQTT/LoRaWAN};
    \draw[very thick, ->, purple!70!black] (2.2, 6.0) -- (2.2, 5.2) node[midway, left] {\scriptsize นโยบายควบคุม $\pi(a|s)$};
    \draw[very thick, ->, teal!70!black] (2.2, 3.0) -- (2.2, 2.2) node[midway, left] {\scriptsize ปรับแรงดัน/ดิวตีไซเคิล};
\end{tikzpicture}
\caption{สถาปัตยกรรม 3 ระดับชั้นของระบบพลาสมากายภาพไซเบอร์เพื่อการเกษตรแม่นยำ}
\label{fig:cpps_architecture}
\end{figure}

\section{การควบคุมแบบวงรอบปิดด้วยการเรียนรู้แบบเสริมกำลัง}

ความท้าทายที่ยิ่งใหญ่ที่สุดของฟิสิกส์พลาสมาในงานเกษตรคือ ความแปรปรวนของสภาวะแวดล้อม (ความชื้นในอากาศ, อุณหภูมิ, คุณภาพน้ำตั้งต้น) ซึ่งส่งผลให้ปฏิกิริยาเคมีการสังเคราะห์ RONS เกิดความไม่แน่นอน การใช้ตัวควบคุมเชิงเส้นแบบดั้งเดิม เช่น PID Controller มักล้มเหลวเนื่องจากระบบมีพฤติกรรมไม่เชิงเส้นขั้นสูงและมีความหน่วงทางเคมี (Chemical latency)

\subsection{การกำหนดสูตรปัญหาแบบมาร์คอฟ}

เราสามารถแปลงปัญหาการควบคุมพลาสมาให้อยู่ในรูปของ \textbf{กระบวนการตัดสินใจแบบมาร์คอฟ (Markov Decision Process หรือ MDP)} ซึ่งประกอบด้วยทูเพิล $(\mathcal{S}, \mathcal{A}, \mathcal{P}, \mathcal{R}, \gamma)$
\begin{enumerate}
    \item \textbf{เวกเตอร์สถานะ ($\mathbf{s}_t \in \mathcal{S}$)} ประกอบด้วยค่าที่วัดได้จากเซนเซอร์แบบเรียลไทม์
    \begin{equation}
        \mathbf{s}_t = \left[ V_{rms}, I_{rms}, T_g, \text{pH}_t, \text{ORP}_t, \text{EC}_t, T_{ambient}, RH\% \right]^T
    \end{equation}
    \item \textbf{เวกเตอร์การกระทำ ($\mathbf{a}_t \in \mathcal{A}$)} คือค่าพารามิเตอร์ที่เอเจนท์ปรับแต่งไปยังแหล่งจ่ายไฟ
    \begin{equation}
        \mathbf{a}_t = \left[ \Delta V_{peak}, \Delta f, \Delta \text{DutyCycle}, \Delta \Phi_{gas} \right]^T
    \end{equation}
    \item \textbf{ฟังก์ชันการให้รางวัล ($\mathcal{R}(\mathbf{s}_t, \mathbf{a}_t)$)} ออกแบบเพื่อเพิ่มอัตราการสังเคราะห์สารออกฤทธิ์เป้าหมาย ขณะเดียวกันก็ลงโทษการสิ้นเปลืองพลังงานและความร้อนที่เกินขีดจำกัดความปลอดภัย
\end{enumerate}

\subsection{อัลกอริทึม Proximal Policy Optimization}

อัลกอริทึมการเรียนรู้แบบเสริมกำลังระดับแนวหน้า เช่น \textbf{Proximal Policy Optimization (PPO)} ถูกนำมาใช้ในการฝึกสอนเอเจนท์ควบคุมพลาสมา เนื่องจากมีความเสถียรสูงและป้องกันไม่ให้นโยบายการควบคุมเปลี่ยนแปลงอย่างรวดเร็วเกินไปจนทำให้ระบบไฟฟ้าแรงสูงเสียหาย ฟังก์ชันเป้าหมายของ PPO กำหนดโดย
\begin{equation}
    L^{CLIP}(\theta) = \widehat{\mathbb{E}}_t \left[ \min\left( r_t(\theta)\widehat{A}_t, \, \text{clip}(r_t(\theta), 1-\epsilon, 1+\epsilon)\widehat{A}_t \right) \right]
\end{equation}
โดยที่ $r_t(\theta) = \frac{\pi_\theta(\mathbf{a}_t|\mathbf{s}_t)}{\pi_{\theta_{old}}(\mathbf{a}_t|\mathbf{s}_t)}$ คืออัตราส่วนความน่าจะเป็นของนโยบาย และ $\widehat{A}_t$ คือฟังก์ชันความได้เปรียบ (Advantage Function) ผลจากการฝึกสอนทำให้ระบบสามารถรักษาระดับค่า ORP ของน้ำ PAW ให้อยู่ที่ $+600\pm15\text{ mV}$ ได้ตลอดเวลา แม้ความชื้นสัมพัทธ์ในบรรยากาศจะเปลี่ยนแปลงอย่างรวดเร็ว

\section{คอมพิวเตอร์วิทัศน์บนขอบสำหรับการประเมินฟีโนไทป์}

การประเมินผลการตอบสนองของเมล็ดพืชและผลผลิตสดหลังผ่านการกระตุ้นด้วยพลาสมา จำเป็นต้องใช้ระบบตรวจวัดที่รวดเร็ว ไม่ทำลายตัวอย่าง และมีความเป็นกลาง เทคโนโลยี \textbf{Edge Computer Vision} ซึ่งผสานโครงข่ายประสาทเทียมตรวจจับวัตถุ (เช่น YOLOv10) และโมเดลจำแนกภาพแบบ Vision Transformer (ViT) ช่วยให้เกษตรกรสามารถประเมินผลลัพธ์ได้ในระดับวินาที

\begin{enumerate}
    \item \textbf{การตรวจวัดพลังการงอกของเมล็ดพันธุ์ (High-throughput Seed Phenotyping)} กล้องอัจฉริยะที่ติดตั้งเหนือถาดเพาะจะบันทึกภาพและประมวลผลหาความยาวของรากแรกเกิด (Radicle length) และอัตราการโผล่พ้นดินของยอดอ่อน เพื่อคำนวณดัชนีความแข็งแรงของเมล็ด (Seed Vigor Index, SVI) อัตโนมัติ
    \item \textbf{การตรวจจับการเข้าทำลายของเชื้อราหลังการเก็บเกี่ยว} อัลกอริทึมสามารถตรวจจับการเจริญของเส้นใยราแอนแทรคโนส (Colletotrichum) หรือราสีเทา (Botrytis cinerea) บนผิวผลไม้ในระยะเริ่มต้น (Incipient stage) ที่ตามนุษย์ยังมองไม่เห็น เพื่อสั่งการให้หัวพ่นพลาสมาเจ็ทเข้าไปทำการฆ่าเชื้อเฉพาะจุดได้อย่างแม่นยำ
\end{enumerate}

\begin{aibox}{การประยุกต์ใช้ Edge AI สำหรับการเกษตรพลาสมาแม่นยำ}
การประมวลผลภาพถ่ายความละเอียดสูงในแปลงเกษตรจริงไม่สามารถพึ่งพาระบบคลาวด์ได้ตลอดเวลา เนื่องจากข้อจำกัดด้านแบนด์วิดท์ของสัญญาณอินเทอร์เน็ตในพื้นที่ห่างไกล จึงจำเป็นต้องบีบอัดและปรับแต่งโมเดลโครงข่ายประสาทเทียมให้สามารถทำงานบนบอร์ดประมวลผลขนาดจิ๋ว (Edge AI Accelerators เช่น NVIDIA Jetson Orin Nano, Google Coral TPU หรือ Raspberry Pi AI Kit)

\textbf{กระบวนการเพิ่มประสิทธิภาพโมเดล (Model Optimization Pipeline):}
\begin{enumerate}
    \item \textbf{Post-Training Quantization (PTQ):} แปลงค่าน้ำหนักของโมเดลจากจุดทศนิยม 32 บิต (FP32) ให้กลายเป็นจำนวนเต็ม 8 บิต (INT8) ช่วยลดขนาดโมเดลลงถึง 4 เท่า และเพิ่มความเร็วในการประมวลผล 3--5 เท่า
    \item \textbf{Structured Pruning:} ตัดทอนน้ำหนักและการเชื่อมต่อของนิวรอนที่ไม่มีนัยสำคัญออก ช่วยลดความซับซ้อนเชิงคำนวณ (FLOPs)
\end{enumerate}
ผลลัพธ์ทำให้โมเดลตรวจจับเมล็ดงอกสามารถรันที่ความเร็วมากกว่า 45 เฟรมต่อวินาที (FPS) โดยใช้พลังงานไฟฟ้าต่ำกว่า 10 วัตต์ ซึ่งสามารถจ่ายไฟเลี้ยงด้วยแผงโซลาร์เซลล์ขนาดเล็กได้อย่างสมบูรณ์
\end{aibox}

\section{แผนที่นำทางสู่การเกษตรคาร์บอนสุทธิเป็นศูนย์}

กระบวนการฮาเบอร์-บอช (Haber-Bosch Process) ที่ใช้ผลิตปุ๋ยไนโตรเจนเคมีในปัจจุบัน ต้องใช้ความดันสูงถึง $15 - 25\text{ MPa}$ และอุณหภูมิสูงกว่า $400 - 500^\circ\text{C}$ ก่อให้เกิดการปล่อยก๊าซคาร์บอนไดออกไซด์คิดเป็นร้อยละ 1.4 ของการปล่อยก๊าซเรือนกระจกทั่วโลก เทคโนโลยีพลาสมาการเกษตรจึงเป็นคำตอบสำคัญในการขับเคลื่อนสู่เป้าหมาย \textbf{Net-Zero Agriculture}

\begin{table}[H]
\caption{แผนที่นำทางเทคโนโลยีพลาสมาการเกษตรและปัญญาประดิษฐ์ (พ.ศ. 2569--2578)}
\label{tab:plasma_roadmap}
\begin{tabularx}{\textwidth}{p{2.6cm} p{3.2cm} X}
\toprule
\textbf{ระยะเวลา} & \textbf{ระดับเทคโนโลยี} & \textbf{เป้าหมายเชิงประจักษ์และการขยายผล} \\
\midrule
ระยะที่ 1 (พ.ศ. 2569--2571) & ระบบพลาสมาอัตโนมัติประจำถิ่น (Decentralized Units) & ติดตั้งสถานีผลิตปุ๋ย PAW พลังงานแสงอาทิตย์ในกลุ่มวิสาหกิจชุมชน ลดการใช้ปุ๋ยเคมีลงร้อยละ 25 \\
ระยะที่ 2 (พ.ศ. 2572--2574) & พลาสมากายภาพไซเบอร์และฝูงหุ่นยนต์ (Swarm Robotics) & บูรณาการหัวพ่นพลาสมาเจ็ทเข้ากับแขนกลหุ่นยนต์กำจัดวัชพืชและแมลงศัตรูพืชแบบไร้สารเคมี \\
ระยะที่ 3 (พ.ศ. 2575--2578) & ระบบนิเวศการเกษตรพลาสมาคาร์บอนสุทธิเป็นศูนย์ & อ่างเก็บกักคาร์บอนทางดินด้วยการตรึงไนโตรเจนชีวภาพร่วมกับ Biochar กระตุ้นด้วยพลาสมา \\
\bottomrule
\end{tabularx}
\end{table}

\begin{examplebox}{ตัวอย่างการออกแบบฟังก์ชันการให้รางวัลสำหรับเอเจนท์ควบคุมพลาสมา}
\textbf{โจทย์} ในระบบผลิตน้ำกระตุ้นด้วยพลาสมา (PAW) จงออกแบบฟังก์ชันการให้รางวัล $R_t$ สำหรับอัลกอริทึม PPO โดยมีเป้าหมายคือ
\begin{itemize}
    \item รักษาค่าศักย์รีดอกซ์ $\text{ORP}$ ให้อยู่ที่ค่าเป้าหมาย $\text{ORP}^* = 600\text{ mV}$
    \item ควบคุมไม่ให้อุณหภูมิของก๊าซ $T_g$ สูงเกินขีดจำกัดความปลอดภัย $T_{max} = 45^\circ\text{C}$
    \item ลดการใช้กำลังไฟฟ้า $P(t)$ ให้เหลือน้อยที่สุด
\end{itemize}

\textbf{การแก้ปัญหา}
\begin{enumerate}
    \item กำหนดเทอมความแม่นยำของสารออกฤทธิ์ (Target Tracking Reward)
    \begin{equation*}
        R_{track} = -\alpha \left( \frac{\text{ORP}_t - \text{ORP}^*}{\text{ORP}^*} \right)^2
    \end{equation*}
    \item กำหนดเทอมลงโทษเมื่ออุณหภูมิเกินเกณฑ์ความปลอดภัย (Thermal Penalty)
    \begin{equation*}
        R_{temp} = -\beta \cdot \max\left( 0, \, T_g - T_{max} \right)^2
    \end{equation*}
    \item กำหนดเทอมประสิทธิภาพเชิงพลังงาน (Energy Efficiency Penalty)
    \begin{equation*}
        R_{energy} = -\gamma \cdot \frac{P(t)}{P_{max}}
    \end{equation*}
    \item สรุปฟังก์ชันการให้รางวัลรวม
    \begin{equation*}
    \begin{aligned}
        R_{total} &= -\alpha \left( \frac{\text{ORP}_t - 600}{600} \right)^2 \\
        &\quad - \beta \cdot \max(0, T_g - 45)^2 - \gamma \cdot \frac{P(t)}{P_{max}}
    \end{aligned}
    \end{equation*}
\end{enumerate}
\textbf{ความหมายทางกายภาพ} ค่าสัมประสิทธิ์ $\alpha, \beta, \gamma$ เป็นน้ำหนักปรับสมดุล (Hyperparameters) หากตั้งค่า $\beta$ ให้มีขนาดใหญ่มาก เอเจนท์จะเรียนรู้โดยอัตโนมัติว่าจะต้องลดดิวตีไซเคิลลงทันทีที่อุณหภูมิเข้าใกล้ 45 องศาเซลเซียส เพื่อป้องกันไม่ให้เกิดความเสียหายต่อหัวพ่นและเนื้อเยื่อพืช
\end{examplebox}

\clearpage

% =========================================================================
% คำถามทบทวนท้ายบทที่ 11 (ขึ้นหน้าใหม่อย่างสง่างาม)
% =========================================================================

\section*{คำถามทบทวนท้ายบทที่ 11}
\addcontentsline{toc}{section}{คำถามทบทวนท้ายบทที่ 11}

\begin{enumerate}
    \item จงอธิบายโครงสร้างและหน้าที่ของระบบพลาสมากายภาพไซเบอร์ (CPPS) ทั้ง 3 ระดับชั้น และอธิบายว่าการเชื่อมต่อข้อมูลระหว่างชั้นส่งผลให้การบำบัดทางการเกษตรมีประสิทธิภาพเหนือกว่าระบบเดิมอย่างไร
    \item เหตุใดการใช้อัลกอริทึมการเรียนรู้แบบเสริมกำลัง (Reinforcement Learning) เช่น PPO จึงมีความเหมาะสมในการควบคุมสภาวะการผลิตน้ำ PAW มากกว่าการใช้ตัวควบคุมแบบ PID ดั้งเดิม
    \item ในการนำโมเดล Deep Learning คอมพิวเตอร์วิทัศน์ไปติดตั้งบนบอร์ดประมวลผลขนาดจิ๋ว (Edge AI) ในแปลงเกษตร เทคนิค Quantization และ Pruning มีความสำคัญอย่างไรต่อการจัดการพลังงานและเวลาหน่วง
    \item จงวิเคราะห์ความเชื่อมโยงระหว่างการใช้เทคโนโลยีพลาสมาผลิตปุ๋ยไนโตรเจนที่แปลงปลูก (On-farm Fertilizer Synthesis) กับเป้าหมายการปล่อยก๊าซเรือนกระจกสุทธิเป็นศูนย์ (Net-Zero Emissions)
    \item หากท่านต้องออกแบบสถานีพลาสมาการเกษตรพลังงานแสงอาทิตย์สำหรับกลุ่มเกษตรกรผู้ปลูกทุเรียนหรือผลไม้ส่งออกในจังหวัดจันทบุรี ท่านจะเลือกใช้ระบบพลาสมาชนิดใด และจะบูรณาการเซนเซอร์และระบบ AI อย่างไรเพื่อควบคุมคุณภาพผลผลิต
\end{enumerate}

\clearpage
